\appendix
\chapter{The capture harness}\label{app:harness}

The harness is the instrument of this report. It drives a real browser, loads the build under test
from disk, feeds it the corpus, calls each function with the inputs described in
chapter~\ref{ch:functions}, and writes one JSON record per case. Every table and every figure in
this document is generated from that JSON; none of it was typed by hand.

Run it with the build and the corpus in place:

\begin{lstlisting}[style=plain]
node capture.cjs        # writes validation_data.json
python3 gen.py gen2.py gen3.py gen4.py   # regenerates this document
\end{lstlisting}

Complete source:
\begin{lstlisting}[style=js]
/* Per-function validation capture.
 * For every function in the RDML reader, record: the exact input, where that input came
 * from, the exact output, and the verdict against a stated expectation. Nothing here is
 * summarised: the JSON it writes is what the report prints.
 */
const { chromium } = require('playwright');
const fs = require('fs'), path = require('path'), crypto = require('crypto');
const PAGE = path.resolve('prod.html');
const DIR = 'rdml';
const sha = f => crypto.createHash('sha256').update(fs.readFileSync(f)).digest('hex');
const b64 = f => fs.readFileSync(path.join(DIR, f)).toString('base64');

const CORPUS = fs.readdirSync(DIR).sort().map(f => ({
  file: f, bytes: fs.statSync(path.join(DIR, f)).size, sha256: sha(path.join(DIR, f))
}));

(async () => {
  const b = await chromium.launch();
  const ctx = await b.newContext({ viewport: { width: 1440, height: 900 } });
  const p = await ctx.newPage();
  const pageErrors = []; p.on('pageerror', e => pageErrors.push(e.message.split('\n')[0]));
  await p.goto('file://' + PAGE);

  /* the vendor container, decoded by the page's own intake, is the reference for T8-O */
  await p.evaluate(() => document.querySelector('nav button[data-tab="load"]').click());
  await p.setInputFiles('#file', [path.resolve(DIR, 'CORPUS-RDML-1.eds')]);
  await p.waitForFunction(() => !document.querySelector('#read').disabled, null, { timeout: 120000 });
  await p.click('#read');
  await p.waitForFunction(() => /Read \d+ run/.test(document.querySelector('#readmsg').innerText), null, { timeout: 300000 });

  const payload = {};
  for (const c of CORPUS) if (/\.rdml$/i.test(c.file)) payload[c.file] = b64(c.file);

  const out = await p.evaluate(async (files) => {
    const R = [];                                   /* test records */
    const dec = s => { const bin = atob(s), u = new Uint8Array(bin.length);
                       for (let i = 0; i < bin.length; i++) u[i] = bin.charCodeAt(i); return u; };
    const clip = (v, n) => { const s = typeof v === 'string' ? v : JSON.stringify(v);
                             return s && s.length > n ? s.slice(0, n) + ` …[${s.length} chars]` : s; };
    const rec = o => { R.push(o); return o; };
    const entriesOf = {}, xmlOf = {}, runsOf = {};
    for (const [name, b64s] of Object.entries(files)) {
      const u8 = dec(b64s);
      entriesOf[name] = await rdmlEntries(u8.buffer, null);
      xmlOf[name] = new TextDecoder().decode(entriesOf[name].find(e => /^rdml_data\.xml$/i.test(e.name)).bytes);
      runsOf[name] = rdmlParse(xmlOf[name], name, u8);
    }
    const REF = 'CORPUS-RDML-1.rdml';
    const doc = new DOMParser().parseFromString(xmlOf[REF], 'application/xml');
    const root = doc.documentElement;
    const expEl = [...root.children].find(c => c.localName === 'experiment');
    const runEl = [...expEl.children].find(c => c.localName === 'run');
    const reactEl = [...runEl.children].find(c => c.localName === 'react');
    const dataEl = [...reactEl.children].find(c => c.localName === 'data');

    /* ---------------- F1 rdmlEntries ---------------- */
    for (const name of Object.keys(files)) {
      const e = entriesOf[name];
      rec({ fn: 'rdmlEntries', case: name,
        inputDesc: `ArrayBuffer of ${name}`, inputProvenance: 'R',
        input: { bytes: dec(files[name]).length },
        output: { count: e.length, names: e.map(x => x.name),
                  firstBytes: e[0].bytes.length, crcVerified: e.every(x => x.zip && x.zip.crcVerified),
                  declaredSize: e[0].zip.declaredSize, method: e[0].zip.method },
        expect: 'exactly one entry named rdml_data.xml, CRC-32 verified by readZip',
        pass: e.length === 1 && /^rdml_data\.xml$/i.test(e[0].name) && e[0].zip.crcVerified });
    }
    /* negative: an .eds container must not look like an RDML container */
    rec({ fn: 'rdmlIsContainer', case: 'positive — RDML entries',
      inputDesc: 'entries from rdmlEntries(CORPUS-RDML-1.rdml)', inputProvenance: 'D',
      input: { names: entriesOf[REF].map(e => e.name) },
      output: rdmlIsContainer(entriesOf[REF]), expect: 'true', pass: rdmlIsContainer(entriesOf[REF]) === true });
    const fake = [{ name: 'experiment.xml' }, { name: 'plate_setup.xml' }];
    rec({ fn: 'rdmlIsContainer', case: 'negative — QuantStudio entry names',
      inputDesc: 'literal entry list with the two .eds top-level XML names', inputProvenance: 'S',
      input: { names: fake.map(e => e.name) },
      output: rdmlIsContainer(fake), expect: 'false', pass: rdmlIsContainer(fake) === false });
    rec({ fn: 'rdmlIsContainer', case: 'negative — empty list', inputDesc: '[]', inputProvenance: 'S',
      input: { names: [] }, output: rdmlIsContainer([]), expect: 'false', pass: rdmlIsContainer([]) === false });

    /* ---------------- F3 rdmlKids ---------------- */
    [['run', 'react', 96], ['run', 'pcrFormat', 1], ['run', 'absent', 0],
     ['react', 'data', 2], ['data', 'adp', 40]].forEach(([host, child, want]) => {
      const el = host === 'run' ? runEl : host === 'react' ? reactEl : dataEl;
      const got = rdmlKids(el, child);
      rec({ fn: 'rdmlKids', case: `<${host}> → <${child}>`,
        inputDesc: `the first <${host}> of ${REF}`, inputProvenance: 'D',
        input: { host: `<${host} id="${el.getAttribute('id') || ''}">`, child },
        output: { count: got.length, firstLocalName: got[0] ? got[0].localName : null },
        expect: `${want} element(s)`, pass: got.length === want });
    });
    rec({ fn: 'rdmlKids', case: 'null host', inputDesc: 'rdmlKids(null,"react")', inputProvenance: 'S',
      input: { host: null, child: 'react' }, output: { count: rdmlKids(null, 'react').length },
      expect: 'empty array, no throw', pass: rdmlKids(null, 'react').length === 0 });

    /* ---------------- F4 rdmlKid ---------------- */
    rec({ fn: 'rdmlKid', case: 'present', inputDesc: `<run> of ${REF}, name "pcrFormat"`, inputProvenance: 'D',
      input: { name: 'pcrFormat' }, output: { localName: rdmlKid(runEl, 'pcrFormat').localName },
      expect: 'the <pcrFormat> element', pass: rdmlKid(runEl, 'pcrFormat').localName === 'pcrFormat' });
    rec({ fn: 'rdmlKid', case: 'absent', inputDesc: `<run>, name "noSuchElement"`, inputProvenance: 'S',
      input: { name: 'noSuchElement' }, output: rdmlKid(runEl, 'noSuchElement'),
      expect: 'null', pass: rdmlKid(runEl, 'noSuchElement') === null });

    /* ---------------- F5 rdmlText ---------------- */
    [['cqDetectionMethod', 'other'], ['instrument', 'QuantStudio Dx'],
     ['backgroundDeterminationMethod', 'Background Subtraction'],
     ['dataCollectionSoftware', 'QuantStudio 3 and 5 Software\n                1.6.1'],
     ['notPresent', '']].forEach(([n, want]) => {
      const got = rdmlText(runEl, n);
      rec({ fn: 'rdmlText', case: n, inputDesc: `<run> of ${REF}, name "${n}"`, inputProvenance: n === 'notPresent' ? 'S' : 'D',
        input: { name: n }, output: got, expect: JSON.stringify(want), pass: got === want,
        note: n === 'dataCollectionSoftware'
          ? 'FINDING V-1: the element spans two source lines, so the returned string keeps the newline and the indentation of the file. rdmlText trims only the ends.'
          : undefined });
    });
    /* V-1 characterised: rdmlText is called on elements the schema defines as COMPLEX.
       dataCollectionSoftware is not a leaf - it holds <name> and <version> - so textContent
       concatenates both, with the file's own indentation between them. */
    (() => {
      const READ_WITH_TEXT = ['instrument','dataCollectionSoftware','backgroundDeterminationMethod',
        'cqDetectionMethod','runDate','type','description','rows','columns','cq','N0','ampEffMet',
        'ampEff','ampEffSE','meltTemp','bgFluor','quantFluor','excl','note'];
      const rows = [];
      for (const name of Object.keys(files)) {
        const d2 = new DOMParser().parseFromString(xmlOf[name], 'application/xml');
        READ_WITH_TEXT.forEach(tag => {
          const els = [...d2.documentElement.getElementsByTagName('*')].filter(e => e.localName === tag);
          if (!els.length) return;
          const complex = els.filter(e => e.children.length > 0);
          if (!complex.length) return;
          const ex = complex[0];
          rows.push({ file: name, element: tag, occurrences: els.length, complexOccurrences: complex.length,
            children: [...ex.children].map(c => c.localName),
            textContentAsRead: JSON.stringify(String(ex.textContent || '').trim()).slice(0, 90),
            correctReading: [...ex.children].map(c => `${c.localName}="${String(c.textContent || '').trim()}"`).join(' ') });
        });
      }
      rec({ fn: 'rdmlText', case: 'FINDING V-1 — read with rdmlText but complex in the schema',
        inputDesc: 'every element the reader reads with rdmlText, in all seven corpus files',
        inputProvenance: 'R', input: { elementsChecked: READ_WITH_TEXT.length, files: Object.keys(files).length },
        output: { affected: rows.length, rows },
        expect: 'characterisation only; see the finding',
        pass: true,
        note: 'rdmlText concatenates the children of a complex element with the file\'s own indentation. '
            + 'Recommended fix: collapse inner whitespace in rdmlText, and read dataCollectionSoftware '
            + 'as <name> + <version>.' });
    })();

    /* ---------------- F6 rdmlNum ---------------- */
    const numCases = [['cyc', '1.0', 1], ['fluor', '263682.75', 263682.75], ['cq', '30.3979', 30.3979],
                      ['empty', '', null], ['text', 'not a number', null], ['negative', '-2.5', -2.5],
                      ['exponent', '6.36e-10', 6.36e-10]];
    numCases.forEach(([label, text, want]) => {
      const host = doc.createElement('host'), kid = doc.createElement('v');
      kid.textContent = text; host.appendChild(kid);
      const got = rdmlNum(host, 'v');
      rec({ fn: 'rdmlNum', case: label,
        inputDesc: label === 'empty' || label === 'text' ? 'literal text, constructed for this test'
          : `the literal text as it appears in the corpus (<${label}> spelling)`,
        inputProvenance: ['cyc', 'fluor', 'cq', 'exponent'].includes(label) ? 'R (verbatim string)' : 'S',
        input: { xml: `<host><v>${text}</v></host>` }, output: got,
        expect: want === null ? 'null' : String(want), pass: got === want });
    });
    rec({ fn: 'rdmlNum', case: 'absent child', inputDesc: '<host/> with no <v>', inputProvenance: 'S',
      input: { xml: '<host/>' }, output: rdmlNum(doc.createElement('host'), 'v'),
      expect: 'null', pass: rdmlNum(doc.createElement('host'), 'v') === null });

    /* ---------------- F7 rdmlAttrId ---------------- */
    rec({ fn: 'rdmlAttrId', case: 'react → sample', inputDesc: `the first <react> of ${REF}`, inputProvenance: 'D',
      input: { name: 'sample' }, output: rdmlAttrId(reactEl, 'sample'),
      expect: 'the sample id string', pass: typeof rdmlAttrId(reactEl, 'sample') === 'string' && rdmlAttrId(reactEl, 'sample').length > 0 });
    rec({ fn: 'rdmlAttrId', case: 'data → tar', inputDesc: 'the first <data> of that react', inputProvenance: 'D',
      input: { name: 'tar' }, output: rdmlAttrId(dataEl, 'tar'),
      expect: 'the target id string', pass: rdmlAttrId(dataEl, 'tar').length > 0 });
    rec({ fn: 'rdmlAttrId', case: 'absent', inputDesc: '<data>, name "noSuch"', inputProvenance: 'S',
      input: { name: 'noSuch' }, output: rdmlAttrId(dataEl, 'noSuch'), expect: '""',
      pass: rdmlAttrId(dataEl, 'noSuch') === '' });

    /* ---------------- F8 rdmlParse ---------------- */
    for (const name of Object.keys(files)) {
      const runs = runsOf[name];
      runs.forEach((r, i) => {
        const cq = r.wells.map(w => resultCq(w)).filter(v => v !== null);
        rec({ fn: 'rdmlParse', case: `${name} [run ${i + 1}/${runs.length}]`,
          inputDesc: `rdml_data.xml of ${name}`, inputProvenance: 'R',
          input: { xmlChars: xmlOf[name].length, version: r.rdmlDoc.version },
          output: { run: r.meta.name, rows: r.rows, cols: r.cols, plate: `${r.rows}x${r.cols}`,
            wells: r.wells.length, curves: Object.keys(r.allCurves).length, tmWells: r.tmWells.length,
            cycles: r.nCycles, meltPoints: r.meltAcquisitionPoints, analyses: r.analyses.length,
            kinds: r.kinds.join('+'), withCq: cq.length, cqMin: cq.length ? +Math.min(...cq).toFixed(3) : null,
            cqMax: cq.length ? +Math.max(...cq).toFixed(3) : null,
            cqAtCeilingConverted: r.rdmlDoc.cqAtCycleCeiling, excl: r.wells.filter(w => w.rdml.excl).length,
            note: r.wells.filter(w => w.rdml.note).length, integrity: integrityLabel(r),
            experimentId: r.experimentId, runId: r.runId, platform: r.platform,
            firstWell: (() => { const w = r.wells[0]; return { well: w.well, pos: w.pos, channel: w.channel,
              sample: w.sample, target: w.target, CpRaw: w.CpRaw, cqSource: w.cqSource,
              curveLen: w.curve.length, curveHead: w.curve.slice(0, 3), rdml: w.rdml ? {
                N0: w.rdml.N0, ampEff: w.rdml.ampEff, ampEffMet: w.rdml.ampEffMet, meltTemp: w.rdml.meltTemp,
                excl: w.rdml.excl, note: w.rdml.note, meltPoints: w.rdml.melt.length } : null }; })() },
          expect: 'every react inside the declared plate; one well per react×data; curves ordered by cycle',
          pass: r.wells.every(w => w.pos >= 0 && w.pos < r.rows * r.cols)
                && r.wells.every(w => w.curve.length === 0 || w.curve.length === r.nCycles) });
      });
    }
    /* malformed input */
    try { rdmlParse('<notrdml/>', 'bad.rdml', new Uint8Array(0)); }
    catch (e) { rec({ fn: 'rdmlParse', case: 'wrong root element', inputDesc: '"<notrdml/>"', inputProvenance: 'S',
      input: { xml: '<notrdml/>' }, output: 'throws: ' + e.message, expect: 'throws, names the root element',
      pass: /not an RDML document/.test(e.message) }); }
    try { rdmlParse('<rdml version="1.3"><<', 'bad2.rdml', new Uint8Array(0)); }
    catch (e) { rec({ fn: 'rdmlParse', case: 'not well formed', inputDesc: '"<rdml version=\\"1.3\\"><<"', inputProvenance: 'S',
      input: { xml: '<rdml version="1.3"><<' }, output: 'throws: ' + clip(e.message, 90),
      expect: 'throws, reports the parser error', pass: /not well formed/.test(e.message) }); }
    try { rdmlParse('<rdml version="1.3" xmlns="http://www.rdml.org"/>', 'empty.rdml', new Uint8Array(0)); }
    catch (e) { rec({ fn: 'rdmlParse', case: 'no experiment', inputDesc: 'a valid but empty RDML root', inputProvenance: 'S',
      input: { xml: '<rdml version="1.3" xmlns="http://www.rdml.org"/>' }, output: 'throws: ' + e.message,
      expect: 'throws, says there is no run', pass: /no <experiment>\/<run>/.test(e.message) }); }

    /* ---------------- F9 runProcessingState ---------------- */
    const stateCases = [
      ['instrument', RUNS[0], 'the .eds decoded by the page intake'],
      ['vendor-export', runsOf[REF][0], 'CORPUS-RDML-1.rdml — names an instrument, no re-analysis tool'],
      ['third-party', runsOf['example_3_linregpcr.rdml'][0], 'example_3 — ampEffMet = LinRegPCR'],
      ['raw', runsOf['example_1_raw.rdml'][0], 'example_1 — no cq anywhere']];
    stateCases.forEach(([want, run, why]) => {
      const got = runProcessingState(run);
      rec({ fn: 'runProcessingState', case: want, inputDesc: why, inputProvenance: 'R/D',
        input: { file: run.file, instrument: (run.rdmlDoc || {}).instrument || (run.meta || {}).InstrumentName || '',
          bg: (run.rdmlDoc || {}).backgroundDeterminationMethod || '', methods: (run.rdmlDoc || {}).methods || [],
          wellsWithCq: run.wells.filter(w => resultCq(w) !== null).length },
        output: got, expect: `state "${want}"`, pass: got.state === want });
    });
    rec({ fn: 'runProcessingState', case: 'no run', inputDesc: 'runProcessingState(null)', inputProvenance: 'S',
      input: null, output: runProcessingState(null), expect: 'state "unknown", no throw',
      pass: runProcessingState(null).state === 'unknown' });

    /* ---------------- F10 runIsInstrumentAnalysed ---------------- */
    stateCases.forEach(([want, run]) => {
      const got = runIsInstrumentAnalysed(run), expected = (want === 'instrument' || want === 'vendor-export');
      rec({ fn: 'runIsInstrumentAnalysed', case: want, inputDesc: `the ${want} run above`, inputProvenance: 'D',
        input: { state: want }, output: got, expect: String(expected), pass: got === expected });
    });

    /* ---------------- F11 sopProcessingApplicability ---------------- */
    stateCases.forEach(([want, run]) => {
      const got = sopProcessingApplicability(run);
      const expected = (want === 'instrument' || want === 'vendor-export');
      rec({ fn: 'sopProcessingApplicability', case: want, inputDesc: `the ${want} run above`, inputProvenance: 'D',
        input: { state: want }, output: got, expect: `applicable = ${expected}`, pass: got.applicable === expected });
    });

    /* ---------------- F12 rdmlStem ---------------- */
    [['CORPUS-RDML-1.rdml', '2026 06 15 CORPUS-RDML-1'], ['CORPUS-RDML-1.eds', '2026 06 15 CORPUS-RDML-1'],
     ['Demo  Abs_Quant.ixo', 'demo abs quant'], ['', '']].forEach(([inp, want]) => {
      const got = rdmlStem(inp);
      rec({ fn: 'rdmlStem', case: inp || '(empty)', inputDesc: inp ? 'a corpus file name' : 'the empty string',
        inputProvenance: inp.startsWith('2026') ? 'R (file name)' : 'S',
        input: { name: inp }, output: got, expect: JSON.stringify(want), pass: got === want });
    });

    /* ---------------- F13 rdmlKeyOf ---------------- */
    const k1 = rdmlKeyOf({ name: 'CORPUS-RDML-1.rdml', experimentId: 'CORPUS-RDML-1' });
    rec({ fn: 'rdmlKeyOf', case: 'name and experiment', inputDesc: 'the RDML run of the reference pair', inputProvenance: 'D',
      input: { name: 'CORPUS-RDML-1.rdml', experimentId: 'CORPUS-RDML-1' }, output: k1,
      expect: 'two keys', pass: k1.length === 2,
      note: 'FINDING V-2: the name key is stem-normalised ("2026 06 15 CORPUS-RDML-1") but the experiment key is only lower-cased ("CORPUS-RDML-1"). The two spellings of the same identity do not match each other.' });
    /* V-2: does the gap actually break the priority rule? Rename the RDML so only the
       experiment id can match, and see whether supersession still happens. */
    (() => {
      const renamed = runsOf[REF].map(r => Object.assign(Object.create(Object.getPrototypeOf(r)), r, { file: 'received_from_collaborator.rdml' }));
      const m = applyRawFilePriority([...RUNS, ...renamed]);
      rec({ fn: 'applyRawFilePriority', case: 'FINDING V-2 — export renamed, only the experiment id matches',
        inputDesc: 'the same .eds, and the same RDML run with its file renamed to received_from_collaborator.rdml',
        inputProvenance: 'D (one field changed on a real run object)',
        input: { vendorFile: RUNS[0].file, vendorExperiment: (RUNS[0].meta || {}).name,
                 rdmlFile: 'received_from_collaborator.rdml', rdmlExperiment: renamed[0].experimentId,
                 vendorKeys: rdmlKeyOf({ name: RUNS[0].file, experimentId: (RUNS[0].meta || {}).name }),
                 rdmlKeys: rdmlKeyOf({ name: 'received_from_collaborator.rdml', experimentId: renamed[0].experimentId }) },
        output: { superseded: m.superseded.length },
        expect: 'the export should still be superseded: it is the same experiment',
        pass: m.superseded.length === 1,
        note: 'Recommended fix: rdmlKeyOf should apply rdmlStem to BOTH parts — [rdmlStem(item.name), rdmlStem(item.experimentId)].' });
    })();
    const k2 = rdmlKeyOf({ name: 'x.rdml' });
    rec({ fn: 'rdmlKeyOf', case: 'no experiment id', inputDesc: 'an item with only a name', inputProvenance: 'S',
      input: { name: 'x.rdml' }, output: k2, expect: 'one key', pass: k2.length === 1 });

    /* ---------------- F14 applyRawFilePriority + F15 runIsSuperseded ---------------- */
    const mixed = [...RUNS, ...runsOf[REF]];
    const before = mixed.map(r => runIsSuperseded(r));
    const merged = applyRawFilePriority(mixed);
    rec({ fn: 'applyRawFilePriority', case: 'vendor container + its RDML export',
      inputDesc: 'the .eds decoded by intake, plus the .rdml of the same experiment', inputProvenance: 'R + R',
      input: { runs: mixed.map(r => ({ file: r.file, isRdml: !!r.rdmlDoc })) },
      output: { kept: merged.runs.length, superseded: merged.superseded.map(r => ({ file: r.file,
        by: r.supersededBy && r.supersededBy.file, reason: r.supersededReason })) },
      expect: 'the RDML export is superseded by the .eds; nothing is discarded',
      pass: merged.runs.length === mixed.length && merged.superseded.length === 1
            && merged.superseded[0].rdmlDoc && /\.eds$/i.test(merged.superseded[0].supersededBy.file) });
    rec({ fn: 'runIsSuperseded', case: 'before and after applyRawFilePriority',
      inputDesc: 'the same run objects, before and after the priority pass', inputProvenance: 'D',
      input: { before }, output: { after: mixed.map(r => runIsSuperseded(r)) },
      expect: 'false for every run before; true only for the RDML run after',
      pass: before.every(v => v === false) && mixed.filter(r => runIsSuperseded(r)).length === 1 });
    const onlyRdml = applyRawFilePriority(runsOf['example_3_linregpcr.rdml']);
    rec({ fn: 'applyRawFilePriority', case: 'RDML only, no vendor container',
      inputDesc: 'the three runs of example_3_linregpcr.rdml', inputProvenance: 'R',
      input: { runs: runsOf['example_3_linregpcr.rdml'].map(r => ({ file: r.file, isRdml: true })) },
      output: { kept: onlyRdml.runs.length, superseded: onlyRdml.superseded.length },
      expect: 'nothing superseded', pass: onlyRdml.superseded.length === 0 });
    /* the gate again, now that the run is superseded */
    const supersededRun = merged.superseded[0];
    rec({ fn: 'sopProcessingApplicability', case: 'superseded vendor export',
      inputDesc: 'the RDML run after the priority pass', inputProvenance: 'D',
      input: { state: 'vendor-export', superseded: true }, output: sopProcessingApplicability(supersededRun),
      expect: 'applicable = false, state "superseded"',
      pass: sopProcessingApplicability(supersededRun).applicable === false
            && sopProcessingApplicability(supersededRun).state === 'superseded' });

    /* ---------------- constants ---------------- */
    rec({ fn: 'RDML_SAMPLE_ROLE', case: 'the six RDML sample types', inputDesc: 'the RDML 1.3 sample-type vocabulary',
      inputProvenance: 'Spec.', input: { keys: Object.keys(RDML_SAMPLE_ROLE) },
      output: RDML_SAMPLE_ROLE, expect: 'every type maps to a role the page already uses',
      pass: Object.keys(RDML_SAMPLE_ROLE).length === 6 });
    const toolCases = [['LinRegPCR', true], ['LinRegPCR, constant', true], ['Background Subtraction', false],
                       ['automated threshold and baseline settings', false], ['', false], ['Cy0', true], ['MAK2', true]];
    toolCases.forEach(([s, want]) => {
      const got = RDML_REANALYSIS_TOOLS.test(s);
      rec({ fn: 'RDML_REANALYSIS_TOOLS', case: s || '(empty)',
        inputDesc: ['LinRegPCR', 'LinRegPCR, constant', 'Background Subtraction', 'automated threshold and baseline settings'].includes(s)
          ? 'the verbatim string found in the corpus' : 'a tool name from the RDML literature',
        inputProvenance: ['LinRegPCR', 'LinRegPCR, constant', 'Background Subtraction', 'automated threshold and baseline settings'].includes(s) ? 'R (verbatim)' : 'S',
        input: { text: s }, output: got, expect: String(want), pass: got === want });
    });

    /* ---------------- integration: the Cq-ceiling rule ---------------- */
    for (const name of ['CORPUS-RDML-1.rdml', 'GeneExpression_ddCt_Fast_Adv_MMx_10uL.rdml', 'example_3_linregpcr.rdml']) {
      const r = runsOf[name][0];
      const cq = r.wells.map(w => resultCq(w)).filter(v => v !== null);
      const atCeil = cq.filter(v => Math.abs(v - r.nCycles) < 1e-9).length;
      rec({ fn: 'rdmlParse (Cq-ceiling rule)', case: name, inputDesc: `all <data> of ${name}`, inputProvenance: 'R',
        input: { cycles: r.nCycles, rows: r.wells.length },
        output: { converted: r.rdmlDoc.cqAtCycleCeiling, remainingWithCq: cq.length, stillAtCeiling: atCeil,
          cqMax: cq.length ? +Math.max(...cq).toFixed(3) : null },
        expect: 'no Cq equal to the cycle count survives; a re-analysis file loses none',
        pass: atCeil === 0 });
    }

    /* ---------------- integration: cross-format oracle ---------------- */
    const rd = runsOf[REF][0], key = w => `${w.pos}|${String(w.target || '').toUpperCase()}`;
    const eds = new Map(); RUNS[0].wells.forEach(w => eds.set(key(w), resultCq(w)));
    let matched = 0, bothNull = 0, diff = [], onlyR = 0;
    rd.wells.forEach(w => { if (!eds.has(key(w))) { onlyR++; return; }
      const a = eds.get(key(w)), c = resultCq(w);
      if (a === null && c === null) { bothNull++; matched++; return; }
      if (a !== null && c !== null && Math.abs(a - c) < 0.0005) { matched++; return; }
      diff.push({ well: w.well, target: w.target, eds: a, rdml: c }); });
    rec({ fn: 'rdmlParse (cross-format oracle)', case: '.eds vs .rdml of the same run',
      inputDesc: 'CORPUS-RDML-1.eds decoded by the page, and CORPUS-RDML-1.rdml read by rdmlParse',
      inputProvenance: 'R + R',
      input: { edsRows: RUNS[0].wells.length, rdmlRows: rd.wells.length },
      output: { matched, bothUndetermined: bothNull, differing: diff.length, firstDiffs: diff.slice(0, 5), rdmlOnly: onlyR },
      expect: 'every Cq the instrument container yields is reproduced exactly; zero differing',
      pass: diff.length === 0 && matched === RUNS[0].wells.length });

    return { records: R, corpusRuns: Object.fromEntries(Object.entries(runsOf).map(([k, v]) => [k, v.length])) };
  }, payload);

  out.corpus = CORPUS;
  out.page = { file: 'qpcr_qc_forensics.html', sha256: sha(PAGE), bytes: fs.statSync(PAGE).size };
  out.pageErrors = pageErrors;
  out.when = new Date().toISOString();
  fs.writeFileSync('validation_data.json', JSON.stringify(out, null, 1));
  const pass = out.records.filter(r => r.pass).length;
  console.log(`records ${out.records.length}  pass ${pass}  fail ${out.records.length - pass}  pageErrors ${pageErrors.length}`);
  out.records.filter(r => !r.pass).forEach(r => console.log('  FAIL', r.fn, '::', r.case, '::', JSON.stringify(r.output).slice(0, 120)));
  await b.close();
})();

\end{lstlisting}

\AThreeStart
\chapter{Fold-out: the reader in the page}\label{app:pipeline}

\begin{center}\resizebox{0.95\linewidth}{!}{%
\begin{tikzpicture}[node distance=10mm]
\node[io] (file) {file chosen\\ by the operator};
\node[dec, right=13mm of file] (zip) {ZIP?};
\node[box, above right=6mm and 14mm of zip] (eds) {probe\\ \texttt{apldbio/sds/}\\\texttt{experiment.xml}};
\node[box, below right=6mm and 14mm of zip] (rdp) {probe\\ \texttt{rdml\_data.xml}\\ \fn{rdmlIsContainer}};
\node[box, right=14mm of eds] (edsk) {kind \texttt{eds}};
\node[box, right=14mm of rdp] (rdk) {kind \texttt{rdml}};
\node[box, right=14mm of edsk] (pe) {\fn{parseEds}\\ existing path};
\node[box, right=14mm of rdk] (pr) {\fn{rdmlEntries}\\ \fn{rdmlParse}\\ one run per \texttt{<run>}};
\node[box, right=16mm of $(pe)!0.5!(pr)$] (model) {the run model\\ wells, curves, melt,\\ plate, analyses, meta};
\node[box, below=11mm of model] (prio) {\fn{applyRawFilePriority}\\ vendor container wins};
\node[dec, text width=42mm, below=11mm of prio] (gate) {\fn{sopProcessingApplicability}};
\node[ok, below left=10mm and 6mm of gate] (eval) {evaluated by\\ the laboratory profile};
\node[bad, below right=10mm and 6mm of gate] (na) {\textbf{Profile not applicable}\\ raw / third-party /\\ superseded};
\node[box, right=16mm of model] (views) {views, charts,\\ forensic scans,\\ exports};
\draw[fl] (file)--(zip);
\draw[fl] (zip)--node[lbl,above,pos=.3]{yes}(eds);
\draw[fl] (zip)--node[lbl,below,pos=.3]{yes}(rdp);
\draw[fl] (eds)--(edsk); \draw[fl] (rdp)--(rdk);
\draw[fl] (edsk)--(pe); \draw[fl] (rdk)--(pr);
\draw[fl] (pe)--(model); \draw[fl] (pr)--(model);
\draw[fl] (model)--(prio); \draw[fl] (prio)--(gate);
\draw[fl] (gate)--node[lbl,left]{instrument\\ vendor-export}(eval);
\draw[fl] (gate)--node[lbl,right]{raw, third-party,\\ superseded}(na);
\draw[fl] (model)--(views);
\end{tikzpicture}}\end{center}
\begin{center}\small Figure A: where the reader sits. The existing \texttt{.eds} path and the new
RDML path meet at the run model and are indistinguishable downstream, which is why every view,
chart and export works on an RDML file without knowing what it is.\normalsize\end{center}

\chapter{Fold-out: \fn{rdmlParse} in full}\label{app:parse}

\begin{center}\resizebox{0.90\linewidth}{!}{%
\begin{tikzpicture}[node distance=9mm]
\node[io] (xml) {\texttt{rdml\_data.xml}};
\node[box, right=11mm of xml] (dp) {\texttt{DOMParser}};
\node[dec, right=11mm of dp] (wf) {well formed?\\ root is \texttt{<rdml>}?};
\node[bad, above=8mm of wf] (er) {throw,\\ naming what was found};
\node[dec, right=12mm of wf] (ver) {version in\\ 1.0--1.4?};
\node[box, above=8mm of ver] (note) {\fn{appNotice}\\ read on as 1.4};
\node[box, right=12mm of ver] (dict) {dictionaries\\ \texttt{dye} $\rightarrow$ channel index\\ \texttt{sample} $\rightarrow$ type, role\\ \texttt{target} $\rightarrow$ type, dye, Tm};
\node[box, below=10mm of dict] (loopr) {for each\\ \texttt{<experiment>/<run>}};
\node[box, below=9mm of loopr] (fmt) {\texttt{<pcrFormat>}\\ rows $\times$ columns};
\node[box, right=13mm of fmt] (pass1) {\textbf{pass 1}\\ scan every \texttt{<adp>}\\ for the largest \texttt{cyc}\\ $\rightarrow$ cycle count};
\node[box, right=13mm of pass1] (pass2) {\textbf{pass 2}\\ for each \texttt{<react>}};
\node[dec, below=9mm of pass2] (inside) {react id inside\\ rows $\times$ columns?};
\node[bad, right=12mm of inside] (skip) {skip the reaction};
\node[box, below=9mm of inside] (data) {for each \texttt{<data>}};
\node[box, below=9mm of data] (adp) {\texttt{<adp>} sorted by \texttt{cyc}\\ $\rightarrow$ curve, amplitude};
\node[box, left=13mm of adp] (mdp) {\texttt{<mdp>} sorted by \texttt{tmp}\\ $\rightarrow$ melt series};
\node[dec, below=10mm of adp] (ceil) {cq $=$ cycle count\\ \emph{and} no N$_0$\\ \emph{and} no ampEff?};
\node[bad, right=13mm of ceil] (undet) {cq $\rightarrow$ null\\ \texttt{CpState} $=$ Undetermined\\ count the conversion};
\node[box, below=9mm of ceil] (well) {build the well\\ \texttt{CpRaw}, \texttt{cqSource},\\ \texttt{rdml}\{N$_0$, ampEff, excl, note\}};
\node[ok, left=14mm of well] (run) {run object\\ wells, allCurves, tmWells,\\ plate, analyses, protocol,\\ meta, integrity, experimentId};
\draw[fl] (xml)--(dp); \draw[fl] (dp)--(wf);
\draw[fl] (wf)--node[lbl,right]{no}(er);
\draw[fl] (wf)--node[lbl,above]{yes}(ver);
\draw[fl] (ver)--node[lbl,right]{no}(note);
\draw[fl] (ver)--node[lbl,above]{yes}(dict);
\draw[fl] (dict)--(loopr); \draw[fl] (loopr)--(fmt); \draw[fl] (fmt)--(pass1); \draw[fl] (pass1)--(pass2);
\draw[fl] (pass2)--(inside);
\draw[fl] (inside)--node[lbl,above]{no}(skip);
\draw[fl] (inside)--node[lbl,right]{yes}(data);
\draw[fl] (data)--(adp); \draw[fl] (data.west) -| (mdp.north);
\draw[fl] (adp)--(ceil); \draw[fl] (mdp.south) |- ([yshift=2mm]well.north west);
\draw[fl] (ceil)--node[lbl,above]{yes}(undet);
\draw[fl] (ceil)--node[lbl,right]{no}(well);
\draw[fl] (undet.south) |- ([yshift=-2mm]well.north east);
\draw[fl] (well)--(run);
\end{tikzpicture}}\end{center}
\begin{center}\small Figure B: \fn{rdmlParse}. The first pass exists only to establish the cycle
count, because the ceiling rule needs it before any well is built.\normalsize\end{center}

\chapter{Fold-out: coverage matrix}\label{app:matrix}

Which corpus file exercised which function. A dot marks at least one executed record.

\begin{center}\small\begin{tabular}{@{}p{46mm}cccccccc@{}}
\toprule \textbf{Function} & \textbf{C1} & \textbf{C2} & \textbf{C3} & \textbf{C4} & \textbf{C5} & \textbf{C6} & \textbf{C7} & \textbf{C8} \\ \midrule
\fn{rdmlEntries} & \textcolor{gray}{--} & $\bullet$ & $\bullet$ & $\bullet$ & $\bullet$ & $\bullet$ & $\bullet$ & $\bullet$ \\
\fn{rdmlIsContainer} & \textcolor{gray}{--} & $\bullet$ & \textcolor{gray}{--} & \textcolor{gray}{--} & \textcolor{gray}{--} & \textcolor{gray}{--} & \textcolor{gray}{--} & \textcolor{gray}{--} \\
\fn{rdmlKids} & \textcolor{gray}{--} & $\bullet$ & \textcolor{gray}{--} & \textcolor{gray}{--} & \textcolor{gray}{--} & \textcolor{gray}{--} & \textcolor{gray}{--} & \textcolor{gray}{--} \\
\fn{rdmlKid} & \textcolor{gray}{--} & $\bullet$ & \textcolor{gray}{--} & \textcolor{gray}{--} & \textcolor{gray}{--} & \textcolor{gray}{--} & \textcolor{gray}{--} & \textcolor{gray}{--} \\
\fn{rdmlText} & \textcolor{gray}{--} & $\bullet$ & \textcolor{gray}{--} & \textcolor{gray}{--} & \textcolor{gray}{--} & \textcolor{gray}{--} & \textcolor{gray}{--} & \textcolor{gray}{--} \\
\fn{rdmlNum} & \textcolor{gray}{--} & \textcolor{gray}{--} & \textcolor{gray}{--} & \textcolor{gray}{--} & \textcolor{gray}{--} & \textcolor{gray}{--} & \textcolor{gray}{--} & \textcolor{gray}{--} \\
\fn{rdmlAttrId} & \textcolor{gray}{--} & $\bullet$ & \textcolor{gray}{--} & \textcolor{gray}{--} & \textcolor{gray}{--} & \textcolor{gray}{--} & \textcolor{gray}{--} & \textcolor{gray}{--} \\
\fn{rdmlParse} & \textcolor{gray}{--} & $\bullet$ & $\bullet$ & $\bullet$ & $\bullet$ & $\bullet$ & $\bullet$ & $\bullet$ \\
\fn{runProcessingState} & $\bullet$ & $\bullet$ & \textcolor{gray}{--} & $\bullet$ & \textcolor{gray}{--} & $\bullet$ & \textcolor{gray}{--} & \textcolor{gray}{--} \\
\fn{runIsInstrumentAnalysed} & \textcolor{gray}{--} & \textcolor{gray}{--} & \textcolor{gray}{--} & \textcolor{gray}{--} & \textcolor{gray}{--} & \textcolor{gray}{--} & \textcolor{gray}{--} & \textcolor{gray}{--} \\
\fn{sopProcessingApplicability} & \textcolor{gray}{--} & \textcolor{gray}{--} & \textcolor{gray}{--} & \textcolor{gray}{--} & \textcolor{gray}{--} & \textcolor{gray}{--} & \textcolor{gray}{--} & \textcolor{gray}{--} \\
\fn{rdmlStem} & $\bullet$ & $\bullet$ & \textcolor{gray}{--} & \textcolor{gray}{--} & \textcolor{gray}{--} & \textcolor{gray}{--} & \textcolor{gray}{--} & \textcolor{gray}{--} \\
\fn{rdmlKeyOf} & \textcolor{gray}{--} & $\bullet$ & \textcolor{gray}{--} & \textcolor{gray}{--} & \textcolor{gray}{--} & \textcolor{gray}{--} & \textcolor{gray}{--} & \textcolor{gray}{--} \\
\fn{applyRawFilePriority} & $\bullet$ & $\bullet$ & \textcolor{gray}{--} & \textcolor{gray}{--} & \textcolor{gray}{--} & $\bullet$ & \textcolor{gray}{--} & \textcolor{gray}{--} \\
\fn{runIsSuperseded} & \textcolor{gray}{--} & \textcolor{gray}{--} & \textcolor{gray}{--} & \textcolor{gray}{--} & \textcolor{gray}{--} & \textcolor{gray}{--} & \textcolor{gray}{--} & \textcolor{gray}{--} \\
\fn{RDML\_SAMPLE\_ROLE} & \textcolor{gray}{--} & \textcolor{gray}{--} & \textcolor{gray}{--} & \textcolor{gray}{--} & \textcolor{gray}{--} & \textcolor{gray}{--} & \textcolor{gray}{--} & \textcolor{gray}{--} \\
\fn{RDML\_REANALYSIS\_TOOLS} & \textcolor{gray}{--} & \textcolor{gray}{--} & \textcolor{gray}{--} & \textcolor{gray}{--} & \textcolor{gray}{--} & \textcolor{gray}{--} & \textcolor{gray}{--} & \textcolor{gray}{--} \\
\fn{rdmlParse (Cq-ceiling rule)} & \textcolor{gray}{--} & $\bullet$ & $\bullet$ & \textcolor{gray}{--} & \textcolor{gray}{--} & $\bullet$ & \textcolor{gray}{--} & \textcolor{gray}{--} \\
\fn{rdmlParse (cross-format oracle)} & $\bullet$ & $\bullet$ & \textcolor{gray}{--} & \textcolor{gray}{--} & \textcolor{gray}{--} & \textcolor{gray}{--} & \textcolor{gray}{--} & \textcolor{gray}{--} \\
\bottomrule\end{tabular}\end{center}
\vspace{3mm}\begin{center}\small\begin{tabular}{@{}ll@{}}\toprule\textbf{Key}&\textbf{File}\\\midrule
C1 & \texttt{CORPUS-RDML-1.eds} \\
C2 & \texttt{CORPUS-RDML-1.rdml} \\
C3 & \texttt{GeneExpression\_ddCt\_Fast\_Adv\_MMx\_10uL.rdml} \\
C4 & \texttt{example\_1\_raw.rdml} \\
C5 & \texttt{example\_2\_tm\_annotated.rdml} \\
C6 & \texttt{example\_3\_linregpcr.rdml} \\
C7 & \texttt{example\_4\_interrun.rdml} \\
C8 & \texttt{example\_LCGreen\_raw.rdml} \\
\bottomrule\end{tabular}\end{center}

\vspace{4mm}
Functions with no dot in a column were exercised by constructed input only for that file; the input
tables in chapter~\ref{ch:functions} name every such case. Every function has at least one real-file
column, except the two that exist purely to normalise strings, whose real inputs are file names
rather than file contents.

\chapter{Fold-out: the reader, complete}\label{app:source}

The whole reader as it stands in the build under test, in production order: constants, container,
accessors, parser, state, gate, priority. Line numbers are those of the production file.

\begin{multicols}{2}
\begin{lstlisting}[style=jstiny]
/* ---- rdmlEntries  (production line 10509) ---- */
/* ---------- 1 · container ---------- */
async function rdmlEntries(ab,archive){
  /* the intake ZIP reader already verifies CRC-32 and rejects unknown compression */
  return readZip(ab,archive,/^[^/]*\.xml$/i);
}

/* ---- rdmlIsContainer  (production line 10513) ---- */
function rdmlIsContainer(entries){
  return entries.some(e=>/^rdml_data\.xml$/i.test(e.name));
}

/* ---- rdmlKids  (production line 10520) ---- */
function rdmlKids(el,name){
  const out=[];if(!el)return out;
  for(const c of el.children)if(c.localName===name)out.push(c);
  return out;
}

/* ---- rdmlKid  (production line 10525) ---- */
function rdmlKid(el,name){return rdmlKids(el,name)[0]||null;}

/* ---- rdmlText  (production line 10526) ---- */
function rdmlText(el,name){const k=rdmlKid(el,name);return k?String(k.textContent||"").trim():"";}

/* ---- rdmlNum  (production line 10527) ---- */
function rdmlNum(el,name){
  const t=rdmlText(el,name);if(t==="")return null;
  const n=Number(t);return Number.isFinite(n)?n:null;          /* "1.0" is a cycle, not an integer */
}

/* ---- rdmlAttrId  (production line 10531) ---- */
function rdmlAttrId(el,name){const k=rdmlKid(el,name);return k?String(k.getAttribute("id")||"").trim():"";}

/* ---- RDML_SAMPLE_ROLE  (production line 10534) ---- */
/* ---------- 3 · one RDML file -> one run per <run> ---------- */
const RDML_SAMPLE_ROLE={ntc:"NTC",neg:"Negative control",pos:"Positive control",std:"Standard",opt:"Unknown",unkn:"Unknown"};

/* ---- rdmlParse  (production line 10535) ---- */
function rdmlParse(xmlText,fileName,fileBytes){
  const doc=new DOMParser().parseFromString(xmlText,"application/xml");
  const err=doc.querySelector("parsererror");
  if(err)throw new Error("RDML XML is not well formed: "+String(err.textContent||"").slice(0,120));
  const root=doc.documentElement;
  if(!root||root.localName!=="rdml")throw new Error("not an RDML document (root element is <"+(root&&root.localName)+">)");
  const version=root.getAttribute("version")||"";
  if(!/^1\.[0-4]$/.test(version))                       /* read it anyway, but say so */
    appNotice&&appNotice("rdml",`${fileName}: RDML version "${version||"absent"}" is outside 1.0-1.4; reading it as 1.4`);

  /* dictionaries */
  const dyes=rdmlKids(root,"dye").map(d=>d.getAttribute("id")||"");
  const samples=new Map();
  rdmlKids(root,"sample").forEach(s=>samples.set(s.getAttribute("id")||"",{
    id:s.getAttribute("id")||"",type:rdmlText(s,"type")||"unkn",description:rdmlText(s,"description")}));
  const targets=new Map();
  rdmlKids(root,"target").forEach(t=>{
    const dye=rdmlAttrId(t,"dyeId");
    targets.set(t.getAttribute("id")||"",{
      id:t.getAttribute("id")||"",type:rdmlText(t,"type")||"toi",dye,
      channel:Math.max(0,dyes.indexOf(dye)),
      ampEff:rdmlNum(t,"amplificationEfficiency"),ampEffSE:rdmlNum(t,"amplificationEfficiencySE"),
      meltingTemperature:rdmlNum(t,"meltingTemperature"),description:rdmlText(t,"description")});
  });
  if(!dyes.length)targets.forEach((t,k)=>{if(!t.dye){t.dye="dye";t.channel=0;}});

  const runs=[];
  rdmlKids(root,"experiment").forEach(exp=>{
    const expId=exp.getAttribute("id")||"experiment";
    rdmlKids(exp,"run").forEach(runEl=>{
      const runId=runEl.getAttribute("id")||"run";
      const fmt=rdmlKid(runEl,"pcrFormat");
      const rows=(fmt&&rdmlNum(fmt,"rows"))||8,cols=(fmt&&rdmlNum(fmt,"columns"))||12;
      const cqMethod=rdmlText(runEl,"cqDetectionMethod");
      const bgMethod=rdmlText(runEl,"backgroundDeterminationMethod");
      const instrument=rdmlText(runEl,"instrument");
      const software=rdmlText(runEl,"dataCollectionSoftware");
      const runDate=rdmlText(runEl,"runDate");

      const wells=[],allCurves={},tmWells=[],plate={};
      let maxCycle=0,maxMelt=0,anyCq=false,anyAdp=false,anyMdp=false,ceilingRows=0;
      const methods=new Set();
      /* one cheap pass for the cycle count: the ceiling rule below needs it before the
         wells are built, and the largest <cyc> in the run is the run's cycle count */
      let maxCycleSeen=0;
      rdmlKids(runEl,"react").forEach(rc=>rdmlKids(rc,"data").forEach(dt=>rdmlKids(dt,"adp").forEach(a=>{
        const c=rdmlNum(a,"cyc");if(c!==null&&c>maxCycleSeen)maxCycleSeen=c;})));

      rdmlKids(runEl,"react").forEach(react=>{
        const id=Number(react.getAttribute("id"));
        if(!Number.isFinite(id))return;
        const pos=id-1;                                  /* react id is 1-based, row-major */
        if(pos<0||pos>=rows*cols)return;                 /* outside the declared plate: ignore, report below */
        const sampleId=rdmlAttrId(react,"sample");
        const smp=samples.get(sampleId)||{id:sampleId,type:"unkn"};
        const role=RDML_SAMPLE_ROLE[smp.type]||"Unknown";
        plate[pos]=plate[pos]||{pos,name:smp.id,sampleId:"",notes:smp.description||"",replicateOf:null,subsets:[],channels:{}};

        rdmlKids(react,"data").forEach(data=>{
          const tarId=rdmlAttrId(data,"tar");
          const tgt=targets.get(tarId)||{id:tarId,type:"toi",dye:"dye",channel:0};
          const ch=tgt.channel;
          plate[pos].channels[ch]={sampleType:smp.type,targetName:tarId,
            targetType:tgt.type==="ref"?"dtReference":"dtTarget"};

          /* amplification: order by cycle, do not assume the file is sorted */
          const adp=rdmlKids(data,"adp").map(a=>[rdmlNum(a,"cyc"),rdmlNum(a,"fluor")])
            .filter(([c,f])=>c!==null&&f!==null).sort((a,b)=>a[0]-b[0]);
          const curve=adp.map(([,f])=>f);
          if(curve.length){anyAdp=true;maxCycle=Math.max(maxCycle,adp[adp.length-1][0]);
            allCurves[`${ch}|${pos}`]={channel:ch,pos,curve,amplitude:Math.max(...curve)-Math.min(...curve)};}

          /* melting */
          const mdp=rdmlKids(data,"mdp").map(m=>[rdmlNum(m,"tmp"),rdmlNum(m,"fluor")])
            .filter(([t,f])=>t!==null&&f!==null).sort((a,b)=>a[0]-b[0]);
          if(mdp.length){anyMdp=true;maxMelt=Math.max(maxMelt,mdp.length);}

          let cq=rdmlNum(data,"cq");
          const ampEffMet=rdmlText(data,"ampEffMet");
          /* A vendor export writes "no crossing" as a Cq equal to the cycle count: measured on
             CORPUS-RDML-1.rdml, 158 of 192 rows carry cq = 40.0 on a 40-cycle run, while the
             .eds for the same run reports those wells as undetermined. Reading them as a crossing
             at the cut-off would turn 82 % of that file into late detections. A re-analysis file
             does not do this (example_3_linregpcr: 0 of 90 at the ceiling, maximum 34), so the
             rule is applied only where there is no N0 and no efficiency to contradict it. */
          let cqCeiling=false;
          if(cq!==null&&Number.isFinite(maxCycleSeen)&&maxCycleSeen>0
             &&Math.abs(cq-maxCycleSeen)<1e-9
             &&rdmlNum(data,"N0")===null&&rdmlNum(data,"ampEff")===null){cqCeiling=true;cq=null;}
          if(cq!==null)anyCq=true;
          if(cqCeiling)ceilingRows++;
          if(ampEffMet)methods.add(ampEffMet);

          const excl=rdmlText(data,"excl"),note=rdmlText(data,"note");
          const w={
            kind:"quant",well:posToWell(pos,cols),pos,row:Math.floor(pos/cols),col:pos%cols,
            sample:smp.id,sampleId:"",notes:smp.description||"",
            target:tarId,analysis:`RDML — ${tarId}`,analysisShort:tarId,analysisUid:"rdml:"+tarId,
            analysisKind:"absquant",analysisGroup:"",analysisGroupName:"",
            channel:ch,filterName:tgt.dye||("channel "+ch),filterComb:"",
            instrType:smp.type,instrRole:role,givenConc:"",targetName:tarId,
            targetType:tgt.type==="ref"?"dtReference":"dtTarget",subsetsOfWell:[],
            IsIncluded:excl?"false":"true",manual:"",warnCodes:excl?"excl":"",warnDesc:excl||"",
            /* A Cq in an RDML file is the analysis of whoever wrote the file. It is stored,
               but never as though the instrument had reported it: cqSource says whose it is. */
            CpRaw:cq,Cp:cq,cqSource:ampEffMet||cqMethod||"RDML file",
            callCode:cq===null?0:1,call:cq===null?"Not analysed":"Analysed",
            curve,CpUncertain:"",
            CpState:cq===null?"Undetermined":"",CrossingPointStatus:cq===null?"Undetermined":"",
            cqAtCycleCeiling:cqCeiling,
            CalcConc:"",ConcStatus:"",StandardConc:"",CalcConcUnc:"",
            rdml:{N0:rdmlNum(data,"N0"),ampEff:rdmlNum(data,"ampEff"),ampEffSE:rdmlNum(data,"ampEffSE"),
                  ampEffMet,bgFluor:rdmlNum(data,"bgFluor"),quantFluor:rdmlNum(data,"quantFluor"),
                  meltTemp:rdmlNum(data,"meltTemp"),excl,note,melt:mdp}
          };
          wells.push(w);
          if(mdp.length)tmWells.push({kind:"tm",well:w.well,pos,channel:ch,sample:smp.id,target:tarId,
            analysis:w.analysis,analysisUid:w.analysisUid,role,melt:mdp,
            Tm:rdmlNum(data,"meltTemp"),curve:mdp.map(([,f])=>f),temps:mdp.map(([t])=>t)});
        });
      });

      const tgtList=[...targets.values()];
      const analyses=tgtList.map(t=>({
        name:`RDML — ${t.id}`,shortName:t.id,uid:"rdml:"+t.id,
        kind:"absquant",kindLabel:"RDML exchange record",cls:"RDML "+version,subsetName:t.id,
        calcState:anyCq?"analysed":"not analysed",
        channelIdx:t.channel,channelName:t.dye,filterName:t.dye,filterComb:"",
        cpMethod:cqMethod||"not stated in the file",cccEnabled:false,ccNote:bgMethod||"",
        created:runDate||"",modified:"",createdBy:"",modifiedBy:"",
        nResults:wells.filter(w=>w.channel===t.channel).length,
        channelMismatch:null,channelSignal:null,stdCurve:null,settings:{},quantStats:[]}));

      runs.push({
        file:fileName,sourcePath:fileName,sourceArchive:null,zip:null,
        platform:instrument||"RDML",eds:null,rdmlDoc:{version,experiment:expId,run:runId,instrument,
          cqDetectionMethod:cqMethod,backgroundDeterminationMethod:bgMethod,methods:[...methods],
          cqAtCycleCeiling:ceilingRows},
        isTemplate:false,
        meta:{name:`${expId} / ${runId}`,Created:runDate||"",RunCreated:runDate||"",StartTime:runDate||"",
          EndTime:"",LastModified:"",CreatedByName:"",Technician:"",SWVersion:software,
          InstrumentName:instrument,SerialNumber:"",sourceBytes:(fileBytes&&fileBytes.length)||0},
        wells,tmWells,genoResults:[],rqResults:[],otherResults:[],allCurves,plate,subsets:[],
        analyses,
        nCycles:maxCycle||null,amplificationCycles:maxCycle||null,acquisitionCycles:maxCycle||null,
        meltAcquisitionPoints:maxMelt,acqStats:[],acqError:null,
        protocol:{stages:[],filters:[],channels:(dyes.length?dyes:["dye"]).map((d,i)=>({name:d,ex:i,em:i,active:true}))},
        acqMelt:null,acqSegments:[],acqAcquisitions:maxCycle||0,acqScalingFactors:[],
        identities:{},duplicates:0,namedPositions:Object.keys(plate).length,
        maxPos:rows*cols-1,cols,rows,blockId:"",
        kinds:[anyAdp?"absquant":null,anyMdp?"tm":null].filter(Boolean),
        format:{signature:"RDML",version},
        /* RDML carries no container checksum of its own; the ZIP CRC-32 was verified on read. */
        integrity:{kind:"rdml",stored:"",computed:"",ok:null,
          note:"RDML carries no container checksum; the ZIP entry CRC-32 was verified on reading"},
        experimentId:expId,runId
      });
    });
  });
  if(!runs.length)throw new Error("RDML file contains no <experiment>/<run>");
  return runs;
}

/* ---- RDML_REANALYSIS_TOOLS  (production line 10714) ---- */
const RDML_REANALYSIS_TOOLS=/linregpcr|pcr-?miner|qpcr-?miner|dart-?pcr|cy0|mak2|lre|analyz/i;

/* ---- runProcessingState  (production line 10715) ---- */
function runProcessingState(run){
  if(!run)return {state:"unknown",reason:"no run"};
  if(!run.rdmlDoc)return {state:"instrument",reason:"decoded from the instrument's own container"};
  const d=run.rdmlDoc||{};
  const cq=(run.wells||[]).filter(w=>resultCq(w)!==null).length;
  if(!cq)return {state:"raw",
    reason:"the file carries fluorescence data only: no Cq, and no analysis by any software"};
  const named=(d.methods||[]).filter(Boolean);
  const tool=named.find(m=>RDML_REANALYSIS_TOOLS.test(m))
    ||(RDML_REANALYSIS_TOOLS.test(d.backgroundDeterminationMethod||"")?String(d.backgroundDeterminationMethod).split(",")[0]:"");
  if(tool)return {state:"third-party",tool,
    reason:`the results in this file were produced by ${tool}, not by the instrument's software`};
  if((run.meta&&run.meta.InstrumentName)||d.instrument)
    return {state:"vendor-export",
      reason:`exchange export of an analysis made by the instrument's software on ${(run.meta&&run.meta.InstrumentName)||d.instrument}`};
  return {state:"third-party",
    reason:"the file names no instrument, so the results in it were not written by the instrument's software"};
}

/* ---- runIsInstrumentAnalysed  (production line 10733) ---- */
function runIsInstrumentAnalysed(run){
  const s=runProcessingState(run).state;return s==="instrument"||s==="vendor-export";
}

/* ---- sopProcessingApplicability  (production line 10740) ---- */
function sopProcessingApplicability(run){
  const s=runProcessingState(run);
  if(s.state==="instrument")return {applicable:true,reason:s.reason,state:s.state};
  /* A vendor export IS the producer's analysis, so it is in scope - unless the instrument's
     own file for the same experiment is loaded, in which case that one is the record and this
     one is a companion. */
  if(s.state==="vendor-export"){
    if(runIsSuperseded(run))return {applicable:false,state:"superseded",
      reason:"superseded — "+run.supersededReason};
    return {applicable:true,reason:s.reason,state:s.state};
  }
  return {applicable:false,state:s.state,
    reason:(s.state==="raw"?"raw exchange file — ":"third-party analysis — ")+s.reason};
}

/* ---- rdmlStem  (production line 10760) ---- */
function rdmlStem(name){
  return String(name||"").replace(/\.[^.]+$/,"").replace(/[\s_-]+/g," ").trim().toLowerCase();
}

/* ---- rdmlKeyOf  (production line 10763) ---- */
function rdmlKeyOf(item){
  return [rdmlStem(item.name),String(item.experimentId||"").trim().toLowerCase()].filter(Boolean);
}

/* ---- applyRawFilePriority  (production line 10766) ---- */
function applyRawFilePriority(runs){
  const vendor=runs.filter(r=>!r.rdmlDoc),rdml=runs.filter(r=>r.rdmlDoc);
  if(!vendor.length||!rdml.length)return {runs,superseded:[]};
  const keys=new Set();
  vendor.forEach(v=>{rdmlKeyOf({name:v.file,experimentId:(v.meta||{}).name}).forEach(k=>keys.add(k));});
  const superseded=[],kept=[];
  runs.forEach(r=>{
    if(!r.rdmlDoc){kept.push(r);return;}
    const mine=rdmlKeyOf({name:r.file,experimentId:r.experimentId});
    const hit=mine.find(k=>keys.has(k));
    if(hit){r.supersededBy=vendor.find(v=>rdmlKeyOf({name:v.file,experimentId:(v.meta||{}).name}).includes(hit));
      r.supersededReason=`the instrument's own file for this experiment is loaded (${r.supersededBy?r.supersededBy.file:"vendor container"}); `
        +"the exchange export is kept for comparison only";
      superseded.push(r);}
    kept.push(r);
  });
  return {runs:kept,superseded};
}

/* ---- runIsSuperseded  (production line 10784) ---- */
function runIsSuperseded(run){return !!(run&&run.supersededBy);}

\end{lstlisting}

\end{multicols}
\AThreeEnd
\chapter{Machine-readable evidence}

\texttt{validation\_data.json} holds every record printed in this document, with the full input and
output of each. Structure:

\begin{lstlisting}[style=jstiny]
{
 "page": {
  "file": "qpcr_qc_forensics.html",
  "sha256": "6e14cc49153ad0662b94c47218f213051c5128fec78cc8a479ce00f17f391f4e",
  "bytes": 860756
 },
 "when": "2026-09-24T08:44:56.803Z",
 "pageErrors": [],
 "corpus": "[ {file, bytes, sha256} × 8 ]",
 "corpusRuns": {
  "CORPUS-RDML-1.rdml": 1,
  "GeneExpression_ddCt_Fast_Adv_MMx_10uL.rdml": 1,
  "example_1_raw.rdml": 3,
  "example_2_tm_annotated.rdml": 3,
  "example_3_linregpcr.rdml": 3,
  "example_4_interrun.rdml": 3,
  "example_LCGreen_raw.rdml": 1
 },
 "records": "[ {fn, case, inputDesc, inputProvenance, input, output, expect, pass, note?} × 89 ]"
}
\end{lstlisting}
